\section{Introduction}

The diffusion of generative AI capabilities since 2022 has reopened a
long-standing question in political sociology: do technological shocks
to the occupational structure erode trust in democratic institutions?
The empirical literature on economic dislocation answers \emph{yes}
within national borders. Trade exposure predicts support for Brexit
\citep{ColantoneStanig2018}; vulnerability to industrial robot adoption
predicts radical-right voting in Western Europe
\citep{AnelliColantoneStanig2021}; and robot installations track voting
shifts in U.S.\ counties \citep{FreyBergerChen2018}. The shared mechanism
is \emph{structural grievance}: technology displaces workers, and
displaced workers express political dissatisfaction
\citep{GidronHall2017,Kurer2020}.

Two features of this evidence base leave the question open for
generative AI. \emph{First}, the dominant exposure measures were built
for manufacturing-heavy automation
\citep{Felten2021,Webb2020,AcemogluRestrepo2020}, but generative AI
exposure concentrates in white-collar and service work
\citep{Gmyrek2025}. The political meaning of dislocation in
administration, retail, and knowledge work need not match that of
robot-driven manufacturing decline. \emph{Second}, almost all of the
existing studies are single-country. That design identifies effects
sharply within a national institutional environment, but it cannot tell
us whether the cross-national pattern matches the within-country one.
\citet{BellJones2015} and \citet{SchmidtCatranFairbrother2016} show
that a multilevel model that constrains the two slopes to be equal can
return a misleading answer when the between- and within-country
associations diverge.

This paper makes the within--between distinction central. It pools the
European Social Survey rounds 6--11 (fielded 2012--2024) and merges
each respondent to the ILO--NASK 2025 GenAI exposure score at 4-digit
ISCO-08 \citep{Gmyrek2025}. From the individual scores it constructs a
country-year mean $\bar E_{ct}$, decomposes that mean into a stable
between-country component $\bar E_c$ and a within-country deviation
$E_{ct}-\bar E_c$, and estimates a sequence of three-level models with
random intercepts at country and country-year.

The headline finding is a sharp gap between the two margins. The
between-country slope is large and positive
($\hat\gamma_B \approx +11.2$, $\mathrm{SE}\approx 3.4$), so countries
with more AI-exposed occupational structures report \emph{higher}
institutional trust. The within-country slope is statistically
indistinguishable from zero, and a Wald test rejects equality.
The education-moderation test of the within-effect also returns a
precise null, offering no support for either a substitution or a curator
reading. The cross-sectional positive correlation therefore appears to
reflect what makes a country high-exposure---income per capita,
service-economy share, institutional capacity---rather than an effect
of AI exposure itself.

Section~\ref{sec:theory} derives the hypotheses from performance,
structural-grievance, and dislocation accounts. Section~\ref{sec:data}
describes the ESS panel, the exposure scores, and the macro covariates.
Section~\ref{sec:methods} sets out the three-level model and the
within--between specification. Section~\ref{sec:results} reports the
M0--M6 progression, the Mundlak Wald, and the robustness battery.
Section~\ref{sec:discussion} interprets the gap and discusses what the
design cannot foreclose.
